% Copy-ready figure blocks. Requires graphicx; compile with XeLaTeX or LuaLaTeX.
% Figure numbering is provisional; figures/ is supplied in prism_figure_reference.zip.
% PDF page 1 / Figure 1
\begin{figure}[p]
\centering
\includegraphics[width=\linewidth,height=0.68\textheight,keepaspectratio]{figures/page_01.pdf}
\caption{\textbf{External docking performance on two supplied-pocket benchmarks.} (A) Astex Diverse Set (85 complexes) and (B) PoseBusters v2 (308 complexes). Solid bar segments indicate selected poses with RMSD $<$2 Å that also pass PoseBusters; hatched extensions indicate RMSD-successful but PB-invalid poses. Total bar length therefore gives RMSD-only success. Local-run results are means over three repeats, with sample SD shown separately at the PB-valid and RMSD-only boundaries. EFF-Dock uses unguided generation with 100 poses and 10 sampling steps, energy-based refinement, and chirality-filtered confidence selection. Other methods retain their evaluated native budgets and selection procedures; SurfDock is represented by its force-optimized variant. GOLD and AutoDock Vina rows below the separator are literature-reported values, not local reruns, and have no repeat SD. Different computational budgets and PoseBusters versions prevent interpreting this as an equal-compute comparison.}
\label{fig:paper-1}
\end{figure}

% PDF page 2 / Figure 2
\begin{figure}[p]
\centering
\includegraphics[width=\linewidth,height=0.68\textheight,keepaspectratio]{figures/page_02.pdf}
\caption{\textbf{Effects of refinement and chirality-aware pose selection.} Selected-pose docking performance is shown for Astex Diverse Set (n=85), PoseBusters v2 (n=308), PhiBench (n=206), FoldBench (n=558), and OpenBind (n=925), using unguided 100-pose, 10-step generation. Raw and energy-refined banks are evaluated with ordinary confidence selection or chirality-filtered confidence selection. Solid segments represent RMSD $<$2 Å and PB-valid success; hatched segments represent RMSD $<$2 Å with PB invalidity. Each total bar height is RMSD-only success. Bars show means over three sampling repeats on the same complexes. Error bars at the solid-segment boundary and total height are the respective sample SDs, not confidence intervals or the SD of the hatched difference. Cohort denominators are retained across conditions.}
\label{fig:paper-2}
\end{figure}

% PDF page 3 / Figure 3
\begin{figure}[p]
\centering
\includegraphics[width=\linewidth,height=0.68\textheight,keepaspectratio]{figures/page_03.pdf}
\caption{\textbf{Docking performance across ligand-complexity strata.} Rows stratify the five external benchmarks by (A) ligand heavy-atom count, (B) RDKit strict rotatable-bond count, and (C) saved rigid-fragment count. Results use unguided 100-pose, 10-step generation, refined candidates, and chirality-filtered confidence selection. Top-1 denotes selected-pose RMSD $<$2 Å success; Top-1 + PB-valid additionally requires the same pose to pass PoseBusters. Oracle denotes at least one RMSD-successful pose among all 100 refined candidates and does not require PB validity. Points and error bars show the mean and sample SD across three repeats. Labels give the number of unique complexes per bin; asterisks identify bins with fewer than five complexes. Empty bins are not imputed. Fragment counts come from saved fragment assignments rather than being inferred from rotatable-bond counts. These strata describe associations with performance, not isolated causal effects of molecular complexity.}
\label{fig:paper-3}
\end{figure}

% PDF page 4 / Figure 4
\begin{figure}[p]
\centering
\includegraphics[width=\linewidth,height=0.68\textheight,keepaspectratio]{figures/page_04.pdf}
\caption{\textbf{Cumulative pose success under confidence ranking and generation order.} (A) The fraction of complexes with at least one RMSD $<$2 Å pose among the k best confidence-ranked candidates from a complete 100-pose bank. (B) The corresponding fraction among the first N candidates in their saved generation order, without confidence sorting. Raw and refined candidates preserve pose identities. Confidence ranking in A is unfiltered, with stable pose-index tie breaking. Lines and shaded bands show means and $\pm$1 sample SD over three repeats for each external benchmark. Solid and dashed lines denote refined and raw banks, respectively. Lower-right insets expand poses 1–20; shaded regions and connectors identify the expanded interval, with the same success-rate scale retained. Both orderings reach the full-bank RMSD oracle at 100 poses. These endpoints measure candidate coverage rather than deployable selected Top-1 success and do not incorporate PoseBusters validity.}
\label{fig:paper-4}
\end{figure}

% PDF page 5 / Figure 5
\begin{figure}[p]
\centering
\includegraphics[width=\linewidth,height=0.68\textheight,keepaspectratio]{figures/page_05.pdf}
\caption{\textbf{Selected-pose success as a function of the available candidate count.} For each of the five external benchmarks, the first N refined poses in saved generation order are retained from the existing unguided 100-pose, 10-step banks. Top-1 reselects the best predicted-RMSD candidate within each prefix; Top-1 + chirality restricts this selection to candidates passing the chirality mask, with fallback to unfiltered Top-1 when no candidate passes. Oracle reports whether the same prefix contains any pose with RMSD $<$2 Å. Lines and shaded bands show the mean and $\pm$1 sample SD across three repeats. All endpoints are RMSD-only because newly selected prefix poses were not subjected to additional official PoseBusters evaluation. Selected Top-1 need not improve monotonically as candidates are added. Reusing saved prefixes is not a separately executed N-pose sampling experiment and does not provide measured runtime at each N.}
\label{fig:paper-5}
\end{figure}

% PDF page 6 / Figure 6
\begin{figure}[p]
\centering
\includegraphics[width=\linewidth,height=0.68\textheight,keepaspectratio]{figures/page_06.pdf}
\caption{\textbf{Guidance and pose–step allocation across post-processing conditions.} Performance on Astex Diverse Set and PoseBusters v2 is compared for 100 poses with 10 sampling steps (N100/S10) and 40 poses with 25 steps (N40/S25), with guidance disabled or enabled at $\eta$=2. The four selection conditions are raw, raw with chirality filtering, refined, and refined with chirality filtering. The endpoint is selected-pose RMSD $<$2 Å together with PoseBusters validity. Points and error bars indicate means and sample SD across three repeats. Guided and unguided runs within a matched budget share the verified initial priors; priors are not assumed nested between the two budgets. Both allocations have the same product N$\times$S, but this does not imply equal wall time because candidate-dependent refinement and confidence evaluation costs differ.}
\label{fig:paper-6}
\end{figure}

% PDF page 7 / Figure 7
\begin{figure}[p]
\centering
\includegraphics[width=\linewidth,height=0.68\textheight,keepaspectratio]{figures/page_07.pdf}
\caption{\textbf{Measured runtime and sampling memory for two unguided budgets.} (A) Aggregate pipeline wall time per complex and (B) peak CUDA allocated memory during sampling for unguided N100/S10 and N40/S25 on Astex Diverse Set and PoseBusters v2. Runtime is the sum of completed shard wall durations divided by the number of complexes within each repeat, averaged across three repeats; whiskers indicate sample SD. It includes recorded pipeline setup, refinement, confidence scoring, and I/O overhead, but excludes separate official PoseBusters evaluation. Because parallel shard durations are summed, this is not user-perceived campaign elapsed time or pure sampling latency. Memory is the maximum recorded sampling allocator peak across shards and repeats, expressed in GiB, and has no SD bar. Measurements use the recorded RTX 6000 Ada runs. Memory is neither whole-pipeline peak usage nor a per-pose quantity.}
\label{fig:paper-7}
\end{figure}

% PDF page 8 / Figure 8
\begin{figure}[p]
\centering
\includegraphics[width=\linewidth,height=0.68\textheight,keepaspectratio]{figures/page_08.pdf}
\caption{\textbf{Sensitivity to pocket crop, prior scale, and pocket-center perturbation.} (A) Sampling-pocket cutoff and (B) initial-prior scale are varied for Astex Diverse Set and PoseBusters v2 in the frozen guided ($\eta$=2) robustness campaign. Rows within each heatmap correspond to Gaussian center-jitter standard deviations of 0, 1, or 2 Å per Cartesian axis. Cells show the mean percentage of selected poses satisfying RMSD $<$2 Å and PoseBusters validity across three repeats, after refinement and ordinary confidence selection without the additional chirality filter. The pocket-cutoff sweep fixes the prior scale at 2 Å, and the prior-scale sweep fixes the sampling cutoff at 10 Å. Only the sampling crop changes in the pocket-cutoff sweep; refinement and scoring retain 10 Å crops. All heatmaps use a shared linear 50–100\% numerical scale, with color contrast emphasized around 70–80\%; the 75\% color anchor is not a statistical threshold. These guided robustness results are distinct from the main unguided evaluation condition.}
\label{fig:paper-8}
\end{figure}

% PDF page 9 / Figure 9
\begin{figure}[p]
\centering
\includegraphics[width=\linewidth,height=0.68\textheight,keepaspectratio]{figures/page_09.pdf}
\caption{\textbf{Sequence relatedness and joint ligand–protein overlap with docking-training data.} All six cohorts are compared with the same 47,277 executed docking-training samples. Eligible ligand-binding chains follow the frozen 6 Å heavy-atom contact definition. (A) Distribution of the maximum binding-chain sequence identity. Here sequence identity is the number of identical known residues, summed over one-to-one binding-chain assignments within a training sample, divided by the total observed query binding-chain length; unaligned residues and unmatched query chains remain in the denominator. (B) Exact canonical heavy-isomeric ligand identity and a sequence score $\geq$70\% are classified by whether they occur in the same training sample, in separate samples only, individually, or neither is observed. All exact-ligand training candidates are searched for the joint condition. (C) The same-sample fraction already contained in B is shown separately with counts and percentages. Cohorts comprise the five external benchmarks and Validation (n=1,076). Sequence bins and the 70\% threshold are descriptive. These results do not establish identical pockets or proven leakage, and no observed match does not certify biological novelty.}
\label{fig:paper-9}
\end{figure}

% PDF page 10 / Figure 10
\begin{figure}[p]
\centering
\includegraphics[width=\linewidth,height=0.68\textheight,keepaspectratio]{figures/page_10.pdf}
\caption{\textbf{Training-ligand relatedness and stratified docking performance.} (A) Composition of each external benchmark by observed exact ligand identity and nearest training-ligand fingerprint similarity. Exact identity uses canonical heavy-atom isomeric SMILES with explicit-hydrogen normalization and takes precedence over similarity bins. Other ligands are grouped by maximum Tanimoto similarity using radius-2, 2,048-bit Morgan fingerprints without chirality. The reference is the parseable ligand index from the 47,277 executed docking-training samples. (B) Selected-pose RMSD $<$2 Å and PB-valid success is compared between complexes with and without an observed exact training-ligand match, using the refined, chirality-filtered unguided condition. Bars and whiskers show the mean and sample SD across three repeats; n denotes complexes per stratum. Ligands are not standardized across tautomers or protonation states. No observed exact match is not proof of complete training-data absence, and group differences are confounded by target and ligand composition rather than being causal estimates of memorization.}
\label{fig:paper-10}
\end{figure}

% PDF page 11 / Figure 11
\begin{figure}[p]
\centering
\includegraphics[width=\linewidth,height=0.68\textheight,keepaspectratio]{figures/page_11.pdf}
\caption{\textbf{Docking performance across training-sequence identity strata.} The five external benchmarks are stratified by maximum binding-chain sequence identity to the 47,277 executed docking-training samples, using bins $<$30\%, 30–$<$70\%, 70–$<$90\%, and 90–100\%. Sequence identity is query-normalized: identical known residues under one-to-one chain assignments are divided by the total observed query binding-chain length, including unaligned residues and unmatched chains. Results use refined candidates and chirality-filtered confidence selection from unguided 100-pose, 10-step banks. Solid segments show RMSD $<$2 Å and PB-valid success; hatched extensions show RMSD-successful but PB-invalid poses. Total bar heights show RMSD-only success. Bars and error bars are means and sample SD across three repeats, with uncertainty shown at both endpoint boundaries. Labels give the number of complexes per bin; empty bins have no estimate. Validation is excluded because no equivalent three-repeat evaluation bank is available. Small strata and target-composition differences limit interpretation of apparent trends.}
\label{fig:paper-11}
\end{figure}

\clearpage
\setcounter{figure}{0}
\renewcommand{\thefigure}{S\arabic{figure}}
% PDF page 12 / Figure S1
\begin{figure}[p]
\centering
\includegraphics[width=\linewidth,height=0.68\textheight,keepaspectratio]{figures/page_12.pdf}
\caption{\textbf{Paired changes in docking success with refinement and chirality filtering.} Success-rate improvements are absolute differences in the percentage of complexes satisfying RMSD $<$2 Å and PoseBusters validity. Panels compare (A) refined versus raw poses without the chirality filter, (B) refined versus raw poses with the filter, (C) filter off versus on for raw poses, and (D) filter off versus on for refined poses. Within each complex, paired differences are averaged across three seeds before resampling. Points show the observed mean improvement in percentage points; whiskers show 95\% percentile confidence intervals from 2,000 bootstrap draws. Complex bootstrap resamples individual complexes. PDB-cluster bootstrap resamples exact-PDB groups while keeping their member complexes together and computes a complex-weighted mean. Both methods share the same point estimate. Exact-PDB grouping does not capture protein-family dependence; PDB-cluster estimates are unavailable for OpenBind. Intervals condition on the three seeds and are not adjusted for multiple comparisons.}
\label{fig:paper-s1}
\end{figure}

% PDF page 13 / Figure S2
\begin{figure}[p]
\centering
\includegraphics[width=\linewidth,height=0.68\textheight,keepaspectratio]{figures/page_13.pdf}
\caption{\textbf{Confidence selection and near-native candidate density.} Refined unguided 100-pose, 10-step banks are evaluated on the five external benchmarks. (A) A complex is counted as successful when at least one pose among the unfiltered confidence-ranked Top-1, Top-5, or all 100 candidates has RMSD $<$2 Å. The Oracle-100 values represent full-bank RMSD coverage. These values summarize the k=1, 5, and 100 positions of the refined confidence-ranked cumulative curves. (B) The percentage of generated candidate poses with RMSD $<$2 Å quantifies near-native candidate density, not geometric diversity. Bars and whiskers show means and sample SD across three repeats. All endpoints are RMSD-only; no full-bank PB-valid fraction or PB-valid Top-5/oracle is inferred. The panels have different units of success: complexes containing a successful candidate in A, versus individual candidate poses in B.}
\label{fig:paper-s2}
\end{figure}

% PDF page 14 / Figure S3
\begin{figure}[p]
\centering
\includegraphics[width=\linewidth,height=0.68\textheight,keepaspectratio]{figures/page_14.pdf}
\caption{\textbf{Descriptive failure decomposition across ligand-complexity strata.} Rows stratify refined, chirality-filtered results by (A) heavy-atom count, (B) strict rotatable-bond count, and (C) saved rigid-fragment count for the five external benchmarks. Each complex belongs to one of four mutually exclusive categories: no refined candidate with RMSD $<$2 Å; an RMSD-successful candidate exists but the selected pose fails the threshold; the selected pose passes RMSD but fails PoseBusters; or the selected pose passes both. Bars show mean fractions over three repeats, using all complexes within each stratum. If O, T, and J denote RMSD oracle coverage, selected Top-1 RMSD success, and selected RMSD-plus-PB success, respectively, the four fractions are 100$-$O, O$-$T, T$-$J, and J. Thus this is a decomposition of the complexity-performance results, not an independent causal experiment. Generation failure refers to the refined bank and does not isolate the sampler from refinement effects; selection failure includes chirality-mask exclusions. Labels give complex counts, and dashes denote empty bins.}
\label{fig:paper-s3}
\end{figure}
